\documentclass[11pt,a4paper]{article}
\usepackage[margin=1in]{geometry}
\usepackage{amsmath,amssymb}
\usepackage{booktabs}
\usepackage{enumitem}
\usepackage{hyperref}
\usepackage{natbib}
\usepackage{tcolorbox}
\usepackage{xcolor}
\usepackage{multirow}
\usepackage{graphicx}

\title{\textbf{Beyond Average ASR: Fine-Grained Vulnerability Analysis and RL-Based Defense for MCP Tool Poisoning}}

\author{
Tianhao Chen \quad Tianchen Guan \quad JiaCheng Sang \\
\textit{Research Proposal --- April 2026}
}
\date{}

\begin{document}
\maketitle

% ============================================================
\begin{abstract}
% ============================================================
Existing evaluations of MCP tool poisoning report aggregate attack success rates (ASR)---e.g., MCPTox's 36.5\% across 20 models---but do not analyze \emph{what kinds of attacks succeed against what kinds of tools on what kinds of models}.
We perform the first fine-grained re-analysis of MCPTox's 25,000+ attack--response records, decomposing ASR along three axes: attack paradigm, tool functional category, and model family.

We report three findings that challenge prevailing assumptions.
\textbf{First}, communication tools (email, messaging) are the most vulnerable category (44.0\% ASR), not read-only tools---contradicting the intuition that low-risk operations are least defended.
\textbf{Second}, attack paradigm (parameter tampering at 50.7\% vs.\ explicit hijacking at 16.7\%) dominates tool category as a predictor of ASR, yet existing defenses do not differentiate by paradigm.
\textbf{Third}, model families cluster distinctly: some models show strong tool-type disparity while others are uniformly vulnerable or uniformly robust, and vulnerability does not monotonically decrease with model scale.

These findings motivate a defense that can discover and adapt to non-obvious vulnerability patterns without manual profiling.
We propose \textbf{MCPDefender}, which uses GRPO (Group Relative Policy Optimization) to train MCP agents that generalize to unseen attack paradigms and tool types.
Unlike DPO, which trains on fixed preference pairs from seen attacks, GRPO's online exploration discovers safety-critical patterns---such as parameter tampering signatures---that static datasets miss.
We evaluate on MCPTox with held-out attack paradigms and tool types, targeting a 10+ percentage point ASR reduction on unseen attacks compared to DPO.
\end{abstract}


% ============================================================
\section{Introduction}
\label{sec:intro}
% ============================================================

MCP tool poisoning is a demonstrated threat.
MCPTox~\citep{mcptox} reports 36.5\% average ASR across 45 real-world MCP servers and 20 LLMs.
MCP-ITP~\citep{mcp-itp} achieves 84.2\% ASR via implicit poisoning.
ToolHijacker~\citep{toolhijacker} reaches 96.7\% through two-phase optimization.

The defense landscape is growing rapidly.
MindGuard~\citep{mindguard} uses attention-based detection, MCPShield~\citep{mcpshield} adds prompt-based security cognition, MCP-Guard trains an E5 embedding classifier, TrustDesc~\citep{trustdesc} generates trusted descriptions from source code, and MCP Safety Training~\citep{mcp-safety-training} applies DPO.
However, the recent MCP-DPT taxonomy~\citep{mcp-dpt} concludes that existing defenses are ``predominantly tool-centric'' with ``persistent gaps,'' and independent evaluations show most defenses succeed less than 30\% of the time.

A fundamental limitation is that both attacks and defenses are evaluated using \textbf{aggregate ASR}---a single number averaging over all tool types, attack paradigms, and models.
This masks critical structure in vulnerability patterns that should inform defense design.

\begin{tcolorbox}[colback=gray!5, colframe=gray!50, boxrule=0.5pt]
\textbf{Motivating question.} If we decompose MCPTox's 36.5\% average ASR by tool category, attack paradigm, and model family, what patterns emerge---and can a defense exploit these patterns automatically?
\end{tcolorbox}

We make three contributions:

\begin{enumerate}[nosep]
    \item \textbf{Fine-grained MCPTox re-analysis} (Section~\ref{sec:analysis}): the first three-way decomposition of MCP poisoning ASR by tool functional category $\times$ attack paradigm $\times$ model family, revealing non-obvious vulnerability patterns that challenge prevailing assumptions.

    \item \textbf{MCPDefender} (Section~\ref{sec:method}): a GRPO-based defense that learns generalizable safety reasoning through online exploration, without requiring manual vulnerability profiling.
    The key advantage over DPO: generalization to \emph{unseen} attack paradigms and tool categories.

    \item \textbf{Systematic defense comparison} (Section~\ref{sec:experiments}): the first apples-to-apples evaluation of SFT, DPO, GRPO, and prompt-based defenses on MCPTox under both seen and unseen attack conditions.
\end{enumerate}


% ============================================================
\section{Related Work}
\label{sec:related}
% ============================================================

\paragraph{MCP attacks.}
MCPTox~\citep{mcptox} defines three paradigms: P1 (explicit hijacking, 36.7\%), P2 (implicit hijacking, 26.7\%), P3 (parameter tampering, 46.7\%).
MCP-ITP~\citep{mcp-itp} automates implicit poisoning.
MPMA~\citep{mpma} demonstrates preference manipulation.
ToolHijacker~\citep{toolhijacker} achieves 96.7\% via two-phase optimization.
All report ASR aggregated over tool types.

\paragraph{MCP defenses.}
MindGuard~\citep{mindguard}: attention-based, white-box only.
MCPShield~\citep{mcpshield}: prompt-based cognitive security.
MCP-Guard: E5 embedding classifier (96\% accuracy).
TrustDesc~\citep{trustdesc}: code-analysis-based description sanitization.
MCP Safety Training~\citep{mcp-safety-training}: DPO on static pairs.
ToolSafe~\citep{toolsafe}: GRPO-trained but for runtime invocation safety, not description poisoning.
TraceSafe evaluates trajectory-level guardrails and finds structural JSON competence correlates more with defense success ($\rho=0.79$) than safety alignment.
No existing defense uses RL to learn generalizable resistance to MCP tool description poisoning.

\paragraph{RL for safety.}
ARLAS~\citep{arlas}: adversarial GRPO for agent safety on BrowserGym/AgentDojo.
MrGuard~\citep{mrguard}: curriculum GRPO for safety classification (+5 F1 over SFT).
RSafe: GRPO-trained judges with superior OOD generalization.
GSPR: generalizable safety policy reasoning via GRPO.
SecAlign~\citep{secalign}: DPO for prompt injection defense.
Our work adapts GRPO specifically to MCP tool poisoning, with evaluation focused on cross-paradigm and cross-tool generalization.


% ============================================================
\section{Fine-Grained MCPTox Re-Analysis}
\label{sec:analysis}
% ============================================================

\subsection{Data and Methodology}

We re-analyze MCPTox's \texttt{response\_all.json}, containing 25,079 attack--response records from 20 victim models across 45 real-world MCP servers.
We annotate each of the 148 mapped tools with a functional category using a 5-level taxonomy based on the principle of least privilege:
L1 (Read-Only, 75 tools), L2 (Create/Modify, 31 tools), L3 (Communicate, 16 tools), L4 (Execute/Privileged, 15 tools), L5 (Destructive, 4 tools).
After filtering unmapped queries and MCPTox-excluded models, 11,203 records are available for analysis.

\subsection{Finding 1: The Inverse Risk Paradox}

\begin{table}[h]
\centering
\small
\caption{Per-category ASR on MCPTox. Communication tools are most vulnerable, not read-only tools.}
\label{tab:level_asr}
\begin{tabular}{lcccc}
\toprule
\textbf{Category} & \textbf{Tools} & \textbf{Trials} & \textbf{ASR} \\
\midrule
L1 Read-Only & 75 & 5,314 & 30.3\% \\
L2 Create/Modify & 31 & 2,475 & 36.0\% \\
L3 Communicate & 16 & 1,183 & \textbf{44.0\%} \\
L4 Execute/Privileged & 15 & 1,345 & 34.0\% \\
L5 Destructive & 4 & 341 & 30.2\% \\
\bottomrule
\end{tabular}
\end{table}

Spearman correlation between ASR and permission level is $\rho = +0.052$ ($p = 0.54$)---no systematic relationship.
The pattern is not risk-monotonic: L3 (communicate) is highest, L1 and L5 are similar.
This contradicts the assumption that low-risk tools are least defended.
We hypothesize that models are trained with explicit ``do not delete'' and ``do not execute arbitrary code'' exemplars, but lack equivalent training for ``do not send data to unauthorized recipients''---making communication tools a blind spot.

\subsection{Finding 2: Attack Paradigm Dominates Tool Category}

\begin{table}[h]
\centering
\small
\caption{Attack paradigm $\times$ tool category ASR. P3 (parameter tampering) is uniformly devastating.}
\label{tab:paradigm_level}
\begin{tabular}{lccc}
\toprule
\textbf{Category} & \textbf{P1 (Explicit)} & \textbf{P2 (Implicit)} & \textbf{P3 (Param.\ Tamper)} \\
\midrule
L1 Read-Only & 22.8\% & 14.9\% & 36.0\% \\
L2 Create/Modify & 17.5\% & 16.4\% & 41.1\% \\
L3 Communicate & 24.2\% & 29.3\% & \textbf{49.4\%} \\
L4 Execute & 16.7\% & 19.9\% & \textbf{50.7\%} \\
L5 Destructive & 2.7\% & 14.2\% & 29.2\% \\
\bottomrule
\end{tabular}
\end{table}

P3 (parameter tampering) is 2--5$\times$ more effective than P1 across all tool categories.
The range across paradigms ($\sim$3\%--51\%) is far larger than across tool categories ($\sim$30\%--44\%).
This means that defense design should prioritize paradigm-specific robustness over tool-type-specific hardening.
Yet no existing MCP defense differentiates by attack paradigm.

\subsection{Finding 3: Model-Family Vulnerability Clustering}

Per-model analysis reveals distinct patterns:
\begin{itemize}[nosep]
    \item Some models show \emph{reversed} disparity: L4 ASR $>$ L1 ASR by 17--20pp (DeepSeek-v3, Qwen3-8b, Gemini-2.5, Mistral).
    \item Some models are near-uniform: $<$5pp gap across all levels (Phi-4, some Qwen variants).
    \item Model scale within families does not consistently reduce vulnerability.
\end{itemize}

These patterns suggest that vulnerability is determined by model-specific training choices, not by general principles---further motivating a learned defense that discovers patterns through exploration rather than relying on prior assumptions.


% ============================================================
\section{Method: MCPDefender}
\label{sec:method}
% ============================================================

\subsection{Motivation: Why DPO is Not Enough}

DPO trains on fixed (chosen, rejected) pairs constructed from known attacks.
If the training data contains only P1 attacks, DPO learns to resist P1 but has no signal for P3's parameter tampering patterns.
Our re-analysis shows P3 accounts for the majority of successful attacks---a defense trained only on P1/P2 pairs will miss the dominant threat.

GRPO addresses this through online exploration.
During training, the agent encounters diverse attack scenarios (including P3-style parameter manipulation) through its own completions.
When the agent fails (selects a malicious tool or constructs tampered arguments), it receives a negative reward and learns from the failure.
This creates training signal for attack patterns that may not be well-represented in static preference datasets.

\subsection{Training Pipeline}

\paragraph{Stage 1: SFT Warm-Start (8 hours).}
Fine-tune Qwen2.5-7B-Instruct with LoRA ($r=64$, 4-bit quantization) on $\sim$3,000 (prompt, safe\_action) pairs constructed from MCPTox's clean tool registries.
Purpose: stabilize JSON output format and establish baseline tool-selection behavior.

\paragraph{Stage 2: DPO Baseline (4 hours).}
Train on preference pairs from MCPTox: chosen = correct tool selection, rejected = malicious tool selection.
This serves both as a baseline and as an initialization for GRPO.

\paragraph{Stage 3: GRPO Training (2--3 days).}
Starting from the DPO checkpoint, run GRPO in the MCP-Gym environment constructed from MCPTox data.

\begin{itemize}[nosep]
    \item \textbf{Environment:} Each episode samples a tool registry (14 clean + 1 poisoned from MCPTox's \texttt{def\_tool/}), a user query, and an attack paradigm.
    50\% of episodes are benign (no poison).
    \item \textbf{Reward:} $R = +1$ (correct tool), $-1$ (hijacked), $-0.5$ (over-refusal on benign), $-0.3$ (format error), $+0.1$ (valid JSON).
    \item \textbf{GRPO config:} $G=8$ completions, $\beta_{\text{KL}}=0.001$, asymmetric clipping, 2,000 steps.
    \item \textbf{Key design choice:} Attack paradigms are sampled uniformly during training.
    GRPO's group-relative advantage naturally allocates more gradient signal to paradigms where the model fails more often (P3 $>$ P1).
\end{itemize}

\subsection{Why GRPO Automatically Discovers Vulnerability Patterns}

Consider a training batch with both P1 and P3 attacks:

For P1 (explicit hijacking, base ASR $\sim$20\%):
Out of $G=8$ completions, $\sim$6 resist correctly ($R=+1$), $\sim$2 fail ($R=-1$).
Group variance is moderate. Advantage signal is moderate.

For P3 (parameter tampering, base ASR $\sim$50\%):
Out of $G=8$ completions, $\sim$4 resist ($R=+1$), $\sim$4 fail ($R=-1$).
Group variance is \textbf{maximal} (binary outcomes, 50/50 split).
The 4 correct completions receive large positive advantages; the 4 failures receive large negative advantages.
\textbf{Gradient signal is strongest exactly where the model needs the most improvement.}

DPO cannot achieve this because it assigns equal weight to all preference pairs regardless of current model vulnerability.
This ``automatic difficulty targeting'' is GRPO's key structural advantage over DPO for MCP defense.


% ============================================================
\section{Experimental Design}
\label{sec:experiments}
% ============================================================

\subsection{Research Questions}

\begin{enumerate}
    \item[\textbf{RQ1.}] What fine-grained vulnerability patterns exist in MCPTox beyond aggregate ASR? (Analysis)
    \item[\textbf{RQ2.}] Does GRPO reduce ASR more than DPO, particularly on unseen attack paradigms? (Core comparison)
    \item[\textbf{RQ3.}] Does GRPO generalize to unseen tool categories? (Transfer)
    \item[\textbf{RQ4.}] What is the source of GRPO's advantage---paradigm targeting, tool-type targeting, or general safety reasoning? (Mechanism)
\end{enumerate}

\subsection{Evaluation: LLM-as-Judge}

A critical methodological point: rule-based judges (keyword matching) fail catastrophically on Template-2 and Template-3 attacks.
Our initial experiments showed a 72 percentage point gap between keyword-judge ASR and MCPTox ground-truth ASR on Template-3 alone.
Following MCPTox's own methodology, we use GPT-4o-mini as an LLM-based judge for all evaluations.
We validate the judge against MCPTox's pre-computed labels, requiring $>$90\% agreement on Template-1 (where keyword judges are reliable) before trusting it on Template-2/3.
At \$0.0001 per call, the total evaluation cost is under \$5.

\subsection{Baselines}

\begin{table}[h]
\centering
\small
\begin{tabular}{llll}
\toprule
\textbf{Method} & \textbf{Type} & \textbf{Online?} & \textbf{Key property} \\
\midrule
No Defense & --- & --- & ASR upper bound \\
Prompt Hardening & Runtime & --- & Safety system prompt \\
DefensiveTokens & Runtime & --- & Prepended defense tokens \\
SFT & Training & Offline & Supervised safe examples \\
DPO & Training & Offline & Static preference pairs \\
\textbf{MCPDefender (GRPO)} & Training & \textbf{Online} & Learned via exploration \\
\bottomrule
\end{tabular}
\caption{Defense baselines. All training-based methods use identical MCPTox data. All evaluations use GPT-4o-mini LLM-as-judge.}
\end{table}

\subsection{Generalization Experiments (Core)}

\paragraph{Cross-paradigm generalization (RQ2).}
The default split is: train on P1 + P2, evaluate on held-out P3.
However, the split is \emph{adaptive}: we first measure the victim model's per-paradigm baseline ASR using LLM-as-judge, then choose the split that maximizes the unseen paradigm's baseline ASR (to ensure sufficient attack surface for demonstrating defense improvements).
\textbf{This is the most important experiment.}
If GRPO generalizes to the unseen paradigm while DPO does not, the paper's main claim is validated.

\paragraph{Cross-tool generalization (RQ3).}
Train on 70\% of MCPTox tools.
Evaluate on 30\% held-out tools.

\paragraph{Cross-model transfer (optional).}
Train on Qwen2.5-7B.
Evaluate zero-shot on Llama-3.1-8B by distilling GRPO's safe traces into DPO data.

\subsection{Expected Results}

\begin{table}[h]
\centering
\small
\caption{Expected main results. The key comparison is \textbf{Unseen P3 ASR}---GRPO should generalize while DPO should not.}
\label{tab:main_results}
\begin{tabular}{l|ccc|cc}
\toprule
\textbf{Method} & \textbf{Seen (P1+P2)$\downarrow$} & \textbf{Unseen (P3)$\downarrow$} & \textbf{All$\downarrow$} & \textbf{BTSR$\uparrow$} & \textbf{ORR$\downarrow$} \\
\midrule
No Defense & $\sim$20 & $\sim$50 & $\sim$32 & $\sim$92 & $\sim$2 \\
Prompt Hardening & $\sim$18 & $\sim$47 & $\sim$30 & $\sim$88 & $\sim$8 \\
SFT & $\sim$14 & $\sim$42 & $\sim$26 & $\sim$90 & $\sim$5 \\
DPO & $\sim$10 & $\sim$\textbf{40} & $\sim$22 & $\sim$88 & $\sim$7 \\
\textbf{MCPDefender} & $\sim$\textbf{8} & $\sim$\textbf{28} & $\sim$\textbf{16} & $\sim$\textbf{89} & $\sim$\textbf{6} \\
\bottomrule
\end{tabular}
\end{table}

\paragraph{Reading the table.}
\begin{itemize}[nosep]
    \item \textbf{Seen (P1+P2):} All training-based methods do well. DPO and GRPO are similar because both saw these patterns during training.
    \item \textbf{Unseen (P3):} DPO barely improves over SFT (40\% vs.\ 42\%) because it never saw P3 patterns. MCPDefender drops to 28\%---a 12pp gap over DPO---because online exploration during GRPO training exposed the model to its own P3-like failures.
    \item \textbf{BTSR/ORR:} MCPDefender maintains utility comparable to SFT.
\end{itemize}

\subsection{Ablation Studies}

\begin{table}[h]
\centering
\small
\begin{tabular}{lll}
\toprule
\textbf{Ablation} & \textbf{Variants} & \textbf{Purpose} \\
\midrule
Training method & SFT / DPO / GRPO & Core comparison \\
GRPO init & From SFT vs.\ from DPO & DPO warm-start value \\
Group size $G$ & 4, 8, 16 & Exploration quality \\
Benign ratio & 30\%, 50\%, 70\% & Over-refusal control \\
Training paradigms & P1 only / P1+P2 / All & Generalization test \\
Training data size & 500 / 1000 / 2000 episodes & Data efficiency \\
\bottomrule
\end{tabular}
\end{table}

\subsection{Mechanism Analysis (RQ4)}

To understand what GRPO learns, we:
\begin{enumerate}[nosep]
    \item Log per-paradigm and per-tool-category advantage magnitudes during training. Expected: P3 advantages are 2--3$\times$ larger than P1 advantages, confirming automatic difficulty targeting.
    \item Extract chain-of-thought from MCPDefender's tool selection reasoning. Analyze whether it develops paradigm-specific detection patterns (e.g., ``this description modifies parameters of an existing tool'' for P3).
    \item Compare GRPO's implicit vulnerability profile (gradient allocation across paradigms/tools) with the ground-truth MCPTox profile from Section~\ref{sec:analysis}.
\end{enumerate}

\subsection{Success Criteria}

\begin{table}[h]
\centering
\small
\begin{tabular}{lcc}
\toprule
\textbf{Criterion} & \textbf{Threshold} & \textbf{Venue} \\
\midrule
MCPDefender Unseen-P3 ASR $<$ DPO Unseen-P3 ASR $-$ 10pp & Required & EMNLP/CCS \\
MCPDefender BTSR $>$ 85\% & Required & Any \\
Per-paradigm advantage plot shows P3 $>$ P1 & Required & EMNLP/CCS \\
MCPDefender unseen-tools ASR $<$ DPO unseen-tools ASR $-$ 5pp & Bonus & Strengthens paper \\
Re-analysis findings are novel and statistically significant & Required & Any \\
\bottomrule
\end{tabular}
\end{table}


% ============================================================
\section{Timeline}
\label{sec:timeline}
% ============================================================

\begin{table}[h]
\centering
\small
\begin{tabular}{lll}
\toprule
\textbf{Phase} & \textbf{Weeks} & \textbf{Deliverable} \\
\midrule
MCPTox re-analysis completion & 1--2 & Findings 1--3, all tables/figures \\
SFT + DPO baselines & 3 & Baseline ASR numbers \\
GRPO training + mechanism logging & 4--5 & MCPDefender checkpoint \\
Generalization + ablation experiments & 6--7 & Main results table \\
DefensiveTokens + prompt hardening baselines & 7 & Complete baseline comparison \\
Analysis + case studies & 8 & Mechanism validation \\
Paper writing & 9--10 & Complete draft \\
\bottomrule
\end{tabular}
\end{table}


% ============================================================
\section{Risks and Mitigations}
% ============================================================

\paragraph{Risk 1: GRPO does not generalize to unseen P3 ($<$5pp gap over DPO).}
Probability: 30\%. Mitigation: The re-analysis findings (Section~\ref{sec:analysis}) stand independently. Reframe as ``analysis + baseline defense comparison'' paper. Target: EMNLP Findings or Workshop.

\paragraph{Risk 2: P3 attacks are too hard for any training-based defense.}
Probability: 15\%. If no method reduces P3 ASR below 40\%, report as a ``hardness result'' that motivates future work. P3's effectiveness itself is a contribution.

\paragraph{Risk 3: MCPTox data quality issues (mapping, labeling).}
Current status: 44.7\% mapping rate. Mitigation: Improve mapping via MCPTox's analysis.ipynb patterns. Supplement with self-run evaluations using MCPTox's \texttt{def\_tool/} on Qwen-7B.

\paragraph{Risk 4: GRPO training instability.}
Mitigation: DPO warm-start provides stable initialization. KL $=0.001$. Curriculum from P1$\to$P2$\to$mixed. Overlong filtering (Bespoke Labs recommendation).


% ============================================================
\section{Expected Contributions}
% ============================================================

\begin{enumerate}
    \item \textbf{First fine-grained MCPTox re-analysis} revealing the inverse risk paradox (communication tools most vulnerable), paradigm dominance (P3 $\gg$ P1), and model-family clustering.

    \item \textbf{MCPDefender}: first RL-based defense for MCP tool poisoning, with demonstrated generalization to unseen attack paradigms.

    \item \textbf{Systematic defense comparison} on MCPTox under both seen and unseen conditions---the first apples-to-apples evaluation of training-based MCP defenses.

    \item \textbf{Mechanism analysis} showing GRPO's automatic difficulty targeting: gradient allocation correlates with vulnerability severity.
\end{enumerate}

\paragraph{Target venues.}
EMNLP 2026 / ACL Rolling (analysis + method), CCS 2026 (if defense results are strong), AISec Workshop (analysis alone suffices).


% ============================================================
\bibliographystyle{plainnat}
\begin{thebibliography}{15}
\bibitem[Wang et~al.(2025a)]{mcptox} Z.~Wang et~al. MCPTox. \emph{AAAI}, 2026.
\bibitem[Li et~al.(2026)]{mcp-itp} R.~Li et~al. MCP-ITP. \emph{arXiv:2601.07395}, 2026.
\bibitem[Wang et~al.(2025b)]{mpma} Z.~Wang et~al. MPMA. \emph{arXiv:2505.11154}, 2025.
\bibitem[Shi et~al.(2025)]{toolhijacker} J.~Shi et~al. ToolHijacker. \emph{NDSS}, 2026.
\bibitem[Wang et~al.(2025c)]{mindguard} Z.~Wang et~al. MindGuard. \emph{arXiv:2508.20412}, 2025.
\bibitem[Zhou et~al.(2026)]{mcpshield} Z.~Zhou et~al. MCPShield. \emph{arXiv:2602.14281}, 2026.
\bibitem[TrustDesc(2026)]{trustdesc} TrustDesc. \emph{arXiv:2604.07536}, 2026.
\bibitem[MCP-DPT(2026)]{mcp-dpt} MCP-DPT. \emph{arXiv:2604.07551}, 2026.
\bibitem[Halloran(2025)]{mcp-safety-training} J.~Halloran. MCP Safety Training. \emph{arXiv:2505.23634}, 2025.
\bibitem[Mou et~al.(2026)]{toolsafe} Y.~Mou et~al. ToolSafe. \emph{arXiv:2601.10156}, 2026.
\bibitem[Chen et~al.(2025)]{secalign} S.~Chen et~al. SecAlign. \emph{CCS}, 2025.
\bibitem[Wang et~al.(2025d)]{arlas} Z.~Wang et~al. ARLAS. \emph{arXiv:2510.05442}, 2025.
\bibitem[Yang et~al.(2025)]{mrguard} MrGuard. \emph{EMNLP}, 2025.
\bibitem[Shao et~al.(2024)]{grpo} Z.~Shao et~al. GRPO. \emph{arXiv:2402.03300}, 2024.
\bibitem[Rafailov et~al.(2023)]{dpo} R.~Rafailov et~al. DPO. \emph{NeurIPS}, 2023.
\end{thebibliography}

\end{document}